\section{Обсуждение}\label{sec:discussion}

\subsection{Интерпретация свода ARC-AGI}

Лучшие по сетке строки для базового обучения (табл.~\ref{tab:arc_base}) показывают, что наибольшая токеновая точность при выборе между режимами ответа \texttt{last} и \texttt{argmax\_interval} достигается у \texttt{mixture\_geometric} при умеренном среднем числе шагов \(\overline{\tau}\approx 4{,}25\). Семейство \texttt{smoothed\_loss} достигает сравнимой точности при \(\overline{\tau}=K\), то есть фактически без экономии глубины в выбранной оптимальной строке сводки. Семейство \texttt{power} даёт конкурентоспособную точность, но с бо́льшими потерями по сравнению с оракульным выбором, что согласуется с менее жёсткой параметризацией хвоста: ошибка распределения времени дороже в смысле~\eqref{eq:oracle_stopping_time}.

\subsection{Дополнительное обучение только головы оракула и сдвиг распределений}

Дополнительное обучение только головы оракула при замороженной рекурсивной части меняет не столько предельную токеновую точность при полном прогоне, сколько совместную конфигурацию «распределение \(\rightarrow\) правило остановки». В табл.~\ref{tab:arc_ft} видно, что для ряда семейств оптимальная стратегия после донастройки сохранённой модели смещается к более ранней остановке при почти неизменной \(\mathrm{tok}_{\mathrm{best}}\). Это можно интерпретировать как смещение соответствия предсказанных масс \(\widehat p_{t,k}\) эмпирическому распределению остаточных ошибок на контрольной выборке: правило \texttt{hazard} становится конкурентоспособным, когда отношение \(\widehat p_{t,k}(k)/\sum_{j\ge k}\widehat p_{t,k}(j)\) лучше согласовано с фактическими условными частотами.

Сдвиг распределения \(\tau^\star\) между обучением и контролем в строгом смысле требует сравнения двух эмпирических распределений на \(\{0,\ldots,K\}\), построенных по независимым наборам задач (обучающий набор и контрольный набор). В данной версии эксперимента параллельная гистограмма \(\tau^\star\) для данных обучения явно не фиксируется; качественное расхождение режимов проявляется в различии кривых токеновой точности по шагам на обучении и контроле и в изменении оптимальной строки сводки после дополнительного обучения оракула. Наблюдаемое для \texttt{negative\_binomial} ухудшение \(\mathrm{tok}_{\mathrm{best}}\) после донастройки согласуется с гипотезой локального переобучения головы на обучающих метках \(\tau^\star\), не совпадающих с режимом сложности на контроле; это \emph{гипотеза}, а не следствие теоремы, и проверяется дополнительными регуляризаторами или ранней остановкой этапа донастройки.

% \subsection{Связь с постановкой математической статистики}

% С позиций теории решений \cite{ferguson1967} оракул~\eqref{eq:oracle_stopping_time} задает нижнюю границу риска при полной информации, а практическое правило~\eqref{eq:real_stopping_rule} --- это \(\Hh_{t,k}\)-измеримая аппроксимация, чей риск раскладывается на ошибку прогноза \(\tau^\star\) и ошибку остановки по глубине. Разделение оценки полного закона и решающей функции позволяет отдельно анализировать чувствительность к сдвигу распределения при смене выборки (например, при переходе от обучения к контролю).
